\begin{textsample}{POS Dim 1 – 2023 – Score 26.00 – 2023/t003333.txt}  \label{ex:f1_pos_063}
Christiana Figueres : There ' s no doubt that we are facing exponential impacts of climate change.
There is also, however, no doubt that that is being met with exponential progress of \textbf{the} technologies that can help to address climate change.
So what we have here is, frankly, a race between two exponential curves.
By 2030, we collectively will have chosen between door number one.
And door number two.
Julio Friedmann : \textbf{The} question I am most commonly asked about climate change is :"Should I be optimistic or pessimistic?"Cynthia Williams : Here ' s what we need to do to make sure that we ' re ready for an all-electric future.
Anika Goss : Being financially secure and climate resilient should be \textbf{the} most important priority.
Faustine Delasalle : It ' s really easy when you look at where we are today and where we ought to be by 2030, to feel like there ' s a huge gap that will never be filled.
But we have good reasons to say that that gap is actually fixable.
Nigel Topping : If you ' re not worried or anxious, you ' re probably not paying attention.
But we should also not be defeatist.
We shouldn ' t say,"We ' ve done nothing, we ' re failing."We ' ve really started 20 years too late, but it takes a long time to get \textbf{the} flywheel of momentum going.
Kingsmill Bond : \textbf{The} tide of change is coming. \textbf{The} costs keep falling. \textbf{The} political pressure keeps increasing.
JF : We ' re going to build a thriving, exciting world.
AG : We can do this.
JF : Full of potential.
AG : We have to do this. [ TED Explores A New Climate Vision ] Manoush Zamorodi : When it comes to climate change, are you an optimist or a pessimist?
News about \textbf{the} climate can be overwhelming and emotional, to \textbf{the} point that many of us just feel numb.
I ' m Manoush Zomorodi, a longtime public media journalist and host of \textbf{the} TED Radio Hour.
And today, we want to bring you something special : an hour of nuance, of real people and clarity about what is happening to our \textbf{planet} and what we ' re doing about it.
Because we all know there ' s plenty of bad news when it comes to \textbf{the} climate crisis.
But guess what?
There ' s also some really good news.
That ' s why we ' re here in Detroit at \textbf{the} TED Countdown Summit, where a cross-section of people who know \textbf{the} most and who are doing \textbf{the} most have gathered to share their ideas and plans with each other and with you.
You ' re going to hear from some incredible speakers about their personal stories and tested solutions.
But first, let ' s get a quick reminder on \textbf{the} basics : climate change in a minute or so. [ Prologue : What is Climate Change? ] David Biello : \textbf{The} air we breathe has changed. \textbf{The} mix of gases is shifting, with more and more \textbf{greenhouse} gases like \textbf{carbon} dioxide.
And \textbf{the} shift is happening faster each year.
In just a few hundred years, \textbf{fossil} fuels formed over eons have been burned as \textbf{coal}, \textbf{oil} and gas. \textbf{The} exhaust has transformed \textbf{the} entire atmosphere and ocean.
It ' s like a pollution blanket.
And \textbf{the} result we know as climate change. \textbf{The} \textbf{planet} has already warmed more than one degree Celsius and is on a path to heat up even more.
That may not seem like a lot, but it ' s already caused major destruction across \textbf{the} globe.
To give us some room to breathe, \textbf{the} world must reduce \textbf{greenhouse} gas pollution by more than seven percent each year, every year of this decade.
There are a number of ways to do that.
Let ' s start by looking at \textbf{the} \textbf{greenhouse} gases themselves.
CO2 makes up nearly 75 percent of \textbf{the} pollution emitted each year.
And then there ' s methane or natural gas, which makes up 17 percent.
Finally, there ' s nitrous oxide, making up six percent of \textbf{the} problem.
There are other \textbf{greenhouse} gases, but these three make up \textbf{the} bulk of \textbf{the} climate challenge.
Where do they come from?
There are several ways to break down \textbf{the} sources of \textbf{greenhouse} gas \textbf{emissions}.
Here we ' re using \textbf{the} public Climate Watch data.
Start with energy. \textbf{The} majority of modern energy comes from burning \textbf{fossil} fuels, which makes \textbf{the} energy sector 76 percent of \textbf{the} climate challenge, including \textbf{the} fuels used for \textbf{transportation}, industrial processes and agricultural production.
Farming and animal raising also have to change, as \textbf{agriculture} accounts for 12 percent of \textbf{greenhouse} gas pollution. \textbf{The} remaining 12 percent comes from a grab bag of human activities like industrial heating, clearing forests and more.
This is \textbf{the} climate challenge we face : going from adding billions of metric tons of \textbf{greenhouse} gases to \textbf{the} air each year to adding zero.
Will it be easy?
No.
Can we do it?
Yes, if we choose to.
MZ : OK, so that ' s \textbf{the} situation.
Where do we go from here?
Well, we ' ve all heard of a tipping point, that moment when an idea or concept finally takes hold and then changes \textbf{the} world.
Well, \textbf{the} UN ' s climate chief, Simon Stiell, believes we are almost at a tipping point for massive action on climate change.
And to get you ready, he wants to take you back to \textbf{the} early ' 90s, when he was an executive at Nokia, and \textbf{the} world wasn ' t sending text messages yet. [ Chapter 1 : Exponential Change ] Simon Stiell : I want you to imagine it ' s 1992, and you ' re talking to a telecoms expert.
She tells you that mobile phone text messaging is going to fundamentally change \textbf{the} way we communicate.
You think, why on Earth am I going to go from this to tapping on a screen with my fingers and my thumbs?
But then, after a slow start, text-based services exploded.
Today, over 23 billion messages are sent every single day.
This paved \textbf{the} way for \textbf{the} \textbf{smartphone} revolution.
I know how much we struggle to comprehend and predict what exponential change means, because I was at Nokia in 1993 when that technological revolution started.
And now, as head of UN climate change, I ' m here to tell you that what was true for mobile phones is also true for climate action today.
We ' re here in Detroit, \textbf{the} heart of car manufacturing.
Take a look at this graph.
Electric vehicle sales are expected to increase to become 70 percent by 2030.
Another example, solar energy.
Between 2010 and 2020, we moved from 20 gigawatts of solar energy installed to 150 gigawatts.
And by 2030, that number is expected to increase again, to 1, 000 gigawatts per year.
Change can come fast.
FD : So exponential change means that change happens very slowly at first and then comes a tipping point.
And suddenly things accelerate and accelerate and accelerate.
NT : It ' s a mathematical term, but it means basically stuff happens very slowly and then it goes really fast.
Tessa Kahn : Governments and experts have consistently underestimated \textbf{the} pace and scale at which we can deploy \textbf{renewable} energy and how quickly \textbf{renewable} energy and associated technologies become affordable.
Ramez Naam : I predicted in 2011 that we ' d have solar as cheap as \textbf{coal} in some parts of \textbf{the} world by 2015 and half \textbf{the} cost of \textbf{coal} by 2020, and everyone thought I was crazy, I ' ve got to say, almost.
And in fact, I was wrong. \textbf{The} actual decline in \textbf{the} cost of solar happened twice as fast as I thought it would.
Solar prices today are a century ahead of where \textbf{the} world ' s leading energy experts thought they ' d be in 2010.
KB : It ' s now spreading from solar, so it ' s happening also in \textbf{the} wind sector, it ' s happening in \textbf{the} battery sector, it ' s happening in \textbf{the} electric vehicle sector.
CF : I don ' t think that we do a good job at communicating \textbf{the} fact that we are progressing more than is commonly known.
Nili Gilbert : \textbf{The} human brain is trained best at thinking about linear change.
But when you think about an exponential curve, next to a line, at first \textbf{the} exponential curve is below \textbf{the} line, and so it ' s moving more slowly than you would think.
And then it crosses \textbf{the} line and it ' s moving faster than you would think.
Habiba Daggash : We call it cautious optimism.
We ' re very hopeful about \textbf{the} trajectories that we ' ve seen.
But we shouldn ' t let up yet.
KB : So we can make this change happen slow or we can make it happen fast. \textbf{The} key point to say is we have agency.
MZ : This idea of exponential growth sounds awesome.
It sounds amazing, but is it actually happening?
Well, I ' m headed into a workshop with industry and climate leaders to find out. \textbf{The} moderators of \textbf{the} workshop challenged us to hold two competing ideas in our minds at \textbf{the} same time.
NT : What we ' re really going to explore today is two apparently opposing realities, and they ' re often two opposing narratives. \textbf{The} one which is true, right, is that in terms of \textbf{the} Paris goals, we ' re way off track.
We ' re deep into \textbf{the} emergency.
But what we ' re going to explore today is \textbf{the} other reality that we know what to do.
Exponential change is happening, but it doesn ' t happen by magic.
MZ : So what are we seeing in different industries?
Well, let ' s start with one that ' s known to be really difficult to decarbonize : cement.
Vicente Saiso Alva : I think everything converges into making climate a huge urgent matter for society.
We started to feel \textbf{the} pressure from all different types of stakeholders, and in one year of running this program, we realized how fast we could move and how fast we were going in reducing CO2 \textbf{emissions}, which means now to go from 40 to close to 50 percent reduction by 2030.
MZ : Encouraging, but to tackle that remaining 50 percent may require technologies that don ' t exist at scale yet.
OK, what about \textbf{shipping}?
Narin Phol : In Maersk, a few years back, we came out with a net-zero ambition by 2050, and in a few years later, based on \textbf{the} data that we have got, we ' ve really seen a path that we can really accelerate that one.
That ' s why \textbf{the} revised target of 2040 came in.
MZ : \textbf{The} news in offshore wind was quite dramatic.
Jennie Dodson : So my moment of mindset shift was in 2019. \textbf{The} government put out a request to bid for offshore wind permits.
And that year \textbf{the} price that came in was two thirds less than just four years before.
And really critically, it was \textbf{the} same price as \textbf{the} mainstream market.
And that was transformative.
Nobody expected that when they started.
MZ : \textbf{The} conversation got pretty technical and a bit wonky.
NP : So to drive exponential growth, you focus on a reinforcing loop to make sure it works and you slow down on \textbf{the} balancing loop.
MZ : But \textbf{the} takeaway is that certain sectors, like EVs and renewables, are being upended, while others, like cement and heavy industry are further behind.
But generally across every industry -- JD : It ' s amazing stuff.
MZ : There ' s an understanding that they all need to work together.
And \textbf{the} message?
Big change can happen.
It ' s going to be hard, but some of us are already doing it, so learn from us. [ Chapter 2 : \textbf{The} Electric Vehicle Revolution ] OK, we are in an exponential mindset.
And as you heard, one example of exponential growth might be sitting in your driveway right now.
Electric vehicles are finally really happening, which means that people here in Detroit are rethinking \textbf{the} auto industry.
They ' re rethinking their jobs, their identity.
Take Cynthia Williams, head of \textbf{sustainability} at Ford.
I sat down with her at Lucky Detroit, a local coffee and barber shop, where she told me that there are three things that need to happen for electric vehicles to go to \textbf{the} next level, from this tipping point to exponential growth.
CW : \textbf{The} three things for me are change, collaboration and capacity.
So \textbf{the} change piece of it is changing \textbf{the} way folks think about \textbf{the} vehicle, making sure that they understand we ' re bringing a better vehicle to \textbf{the} customers.
MZ : So you don ' t have to make a sacrifice.
CW : I think one of \textbf{the} things we need to do is actually get people in \textbf{the} cars and show them how fun it is.
And it ' s no different, it ' s just powered by a battery.
We are at a tipping point.
We ' re at a point of acceleration of electric vehicles.
Many are calling it \textbf{the} electric vehicle revolution.
Now with that comes unprecedented change that will require us to build new capacity and collaborate in ways that we ' ve never seen before.
To make things happen faster, you have to think differently.
You can ' t just collaborate with \textbf{the} same people because you get \textbf{the} same answer, right?
You need to broaden your collaboration space.
Go outside of your industry, understand what worked in Norway, for example, to get them to where they are.
MZ : Where ' s \textbf{the} weak link, Cynthia?
What ' s \textbf{the} thing that you ' re like, that if we don ' t push through that, we may not get to this full adoption?
CW : \textbf{The} critical things that we hear from \textbf{the} consumers is they need infrastructure.
That ' s what it ' s taking for us to work with governments to make sure that we bring more charging to everyone, whether it ' s rural, urban, wealthy, low-income, we all need chargers.
And I ' m not saying we need one on every corner, but we do need many, many more.
When we start building out \textbf{the} charging infrastructure, we need lighting.
We need, you know, security, shelter, we need restrooms.
We need other amenities so that people know where to go, where to charge.
Next comes capacity.
We have to invest in new facilities and talent.
Now, that ' s true for any industry transitioning to greener goals and strategies.
And we ' re investing in talent, we ' re investing in education and training, and we ' re also listening to \textbf{the} communities where we live and operate.
MZ : So a young Cynthia Williams, now, your counterpart who ' s studying engineering, will she be learning something very different if she wants to go and work in \textbf{the} automotive industry in Detroit?
CW : There will be different options, right.
And so, I ' m a mechanical engineer by trade, but I think there will be opportunities for \textbf{electrical} engineering, software engineering.
And there ' s even certificates that students can get instead of going through a four-year college.
It ' s \textbf{the} proper training to get them right, to be dedicated, to build \textbf{the} vehicles that we need for \textbf{the} future.
At \textbf{the} end of this year, globally, we ' ll have \textbf{the} ability to produce 600, 000 electric vehicles.
And in 2026, we ' ll have a global goal to produce two million electric vehicles.
And it just goes up from there.
MZ : What is this going to look like, this growth for me in \textbf{the} next couple of years, in five years, in 20 years?
CW : Clean air.
That ' s what I look at.
If you think about when we were, during \textbf{the} pandemic, when everybody stopped driving and they were at home, clear skies and moving to zero tailpipe \textbf{emission} vehicles.
And if you chose an electric vehicle, that ' s you participating and you contributing to a better world, to a better future.
We have come a long way.
It ' s so gratifying to see electric vehicles on \textbf{the} road, to move from aspiration to reality.
To see plants dedicated to building electric vehicles.
We have \textbf{the} technology, we have \textbf{the} leadership, we have \textbf{the} ingenuity to respond to \textbf{the} environmental crisis with courage.
And I ' ll tell you, as a person who ' s dedicated my entire career to \textbf{sustainability}, \textbf{the} road ahead is paved with promise.
TED-Ed : If you were buying a car in 1899, you would have had three major options to choose from.
You could buy a steam-powered car, typically relying on gas-powered boilers.
These could drive as far as you wanted, provided you also wanted to lug around extra water to refuel and didn ' t mind waiting 30 minutes for your engine to heat up.
Alternatively, you could buy a car powered by gasoline.
However, \textbf{the} internal combustion engines in these models required dangerous hand cranking to start and emitted loud noises and foul smelling exhaust while driving.
So your best bet was probably option number three : a battery-powered electric vehicle.
These cars were quick to start, clean and quiet to run, and if you lived somewhere with access to electricity, easy to refuel overnight.
If this seems like an easy choice, you ' re not alone.
By \textbf{the} end of \textbf{the} 19th century, nearly 40 percent of American cars were electric.
In cities with early electric systems, battery-powered cars were a popular and reliable alternative to their occasionally explosive competitors.
But electric vehicles had one major problem.
Batteries.
MZ : Flash forward to \textbf{the} present and batteries are still on everyone ' s mind.
So I visited Mujeeb Ijaz, \textbf{the} founder of \textbf{the} energy storage company One, who also happens to be an electric vehicle history buff.
Mujeeb Ijaz : So this is a 1922 Detroit Electric.
MZ : OK, so 100-year-old vehicle.
MI : That ' s right, electric vehicle.
And then this was Walter Baker.
He developed \textbf{the} Baker Electric.
And he made \textbf{the} range run 201 miles on a single charge.
He also felt strongly that you should be able to wear your top hat, you know, so he has a taller carriage.
These had lead acid batteries to begin with.
And then Thomas Edison introduced \textbf{the} nickel-iron battery.
So this is \textbf{the} battery compartment, one of \textbf{the} battery compartments.
And this is where you would charge.
MZ : Right there?
MI : And we ' re building on what these original pioneers did.
And I think it ' s important to not think about us starting from scratch now.
But go \textbf{backwards}, think about how they solved \textbf{the} problem, what disrupted their ability to advance this to \textbf{the} market, and then pick those problems up and make sure we don ' t get stuck again.
MZ : So we ' ve been here before, a moment when electric vehicles could have taken off but didn ' t.
So what do we need to do differently this time?
MI : Electric vehicles were able to deal with mobility in urban areas, but there were charging deserts where rural driving resulted in no way to get back because you couldn ' t find electric charging everywhere.
And that was a problem.
MZ : And that ' s what people are worried about right now.
We ' re at \textbf{the} same thing.
We ' re worried about range, we ' re worried about accessibility.
All of \textbf{the} things are \textbf{the} same.
But \textbf{the} key difference being that we now know that our \textbf{planet} can ' t keep going like this.
MI : That ' s right.
Had \textbf{the} vision of \textbf{the} Model T assembly line and mass manufacturing been applied to electric car, \textbf{the} electric car would have also dropped in price, but range and charging deserts would have still been \textbf{the} obstruction.
MZ : OK.
We didn ' t have \textbf{the} technology, but we do now.
MI : We do now.
MZ : In \textbf{the} US, there hasn ' t been a mass move to electric vehicles... yet.
Mujeeb ' s company is one of many who want to change that by building batteries that give drivers longer range.
And these batteries aren ' t just for cars.
MI : So we produce battery packs for \textbf{transportation}, but we ' ve also started developing \textbf{the} same battery products into grid batteries.
So micro-grids are those technologies that combine solar and storage and replace a \textbf{coal} plant or a natural gas plant for generating electricity.
So buildings and industry and \textbf{transportation} all working together.
And that ' s an exciting future, is all these industries working together, that makes me more confident that we ' re at \textbf{the} time in history where this transition will take place.
MZ : How many years till we get to that?
MI : I think it ' s within 10 years.
MZ : No way, 10 years?
MI : We ' re in a place where that transition is going to be much more visible, and then we just need technologies to drive it up to \textbf{the} mass markets.
MZ : I am so excited.
Like, genuinely so excited.
Everyone everywhere driving electric sounds great, but that also means that \textbf{the} power to charge all those vehicles needs to be green.
And that brings us to \textbf{the} energy sector, which is also changing at an incredible pace, even though \textbf{renewable} energy has been years in \textbf{the} making. [ Chapter 3 : Renewables ] TED-Ed : In \textbf{the} \textbf{spring} of 1954, \textbf{the} press excitedly gathered around Bell Laboratories ' latest invention : a silicon-based solar cell that could efficiently convert \textbf{the} sun ' s energy into \textbf{electrical} current. \textbf{The} creation was celebrated as \textbf{the} dawn of a new era, as reporters touted that civilization would soon run on \textbf{the} sun ' s limitless energy.
But \textbf{the} dream had a catch.
As this first commercially sold solar cell cost around 300 dollars per watt, meaning at its current rate it would cost well over a million dollars to buy a unit large enough to power a single home.
But today, in many countries, solar is \textbf{the} cheapest form of energy to produce, surpassing \textbf{fossil} fuel alternatives like \textbf{coal} and natural gas.
Millions of homes are equipped with rooftop solar, with most units paying for themselves in their first seven to 12 years and then generating further savings.
So how did solar become so affordable?
A turning point in solar ' s price history \textbf{occurred} on \textbf{the} floor of Germany ' s parliament, where in 2000, Hermann Scheer introduced \textbf{the} \textbf{Renewable} Energy Sources Act.
This legislation laid out a vision for \textbf{the} country ' s energy future in solar and wind.
It incentivized citizens to personally invest in rooftop solar panels by guaranteeing payment to homeowners for \textbf{the} \textbf{renewable} energy they generated and sold to \textbf{the} grid. \textbf{The} pay rate for this electricity was highly subsidized, at times reaching four times \textbf{the} market price.
Several other countries soon followed Germany ' s example, implementing similar policies and incentives to drive their country ' s solar use.
This created unprecedented demand for solar panels worldwide.
Manufacturers were able to scale up production and innovate in ways that cut costs.
As a result, solar panel prices dropped while efficiency grew.
HD : We see \textbf{the} cost reduction in solar and other \textbf{renewable} energy technologies has been so steep, but a lot of it is because of \textbf{the} enabling regulatory and policy environments created in places like \textbf{the} European Union and \textbf{the} US.
Akil Callender : We cannot divert catastrophic climate change without addressing \textbf{emissions} coming from our energy systems.
Currently, \textbf{emissions} from \textbf{the} energy system account for around two thirds of global \textbf{greenhouse} gas \textbf{emissions}.
So \textbf{the} good news is, yes, we are seeing tremendous progress across \textbf{the} globe in \textbf{renewable} energy.
It ' s rolling out faster than ever before.
Luisa Neubauer : \textbf{The} brilliant thing about energy supply in \textbf{the} world is that we have alternatives.
So we don ' t need \textbf{fossil} fuels to provide safe and equitable access to energy systems.
Ramón Méndez Galain : Almost 90 percent of \textbf{the} total new installed capacity all over \textbf{the} world is just \textbf{renewable}.
This is a tremendous change.
AC : There ' s more finance than ever before going into clean energy.
In many countries, it is \textbf{the} largest contributor to their electricity system for \textbf{the} first time ever.
KB : And also because China is \textbf{the} leading manufacturing nation in \textbf{the} world, change is happening there much faster.
And those innovations are then spreading around \textbf{the} world.
Hongquiao Liu : China ' s \textbf{coal} demand might not have peaked, and China will continue to consume quite a lot of \textbf{coal} in \textbf{the} next few years.
But \textbf{the} decarbonization is happening.
It ' s happening faster than anyone has expected, and China is adding an astonishing amount of wind and solar to its grid every year.
Eduarda Zoghbi : I think one of \textbf{the} biggest obstacles today are natural resources, in \textbf{the} sense that we ' re talking more and more about critical minerals and \textbf{the} role that they have.
And batteries are demanding a lot of critical minerals that are present in most of \textbf{the} developing countries.
So I think that even though costs are decreasing, we have to be more and more creative when it comes to recycling \textbf{the} batteries.
We also have to think of how they ' re impacting local communities.
KB : So to give you \textbf{the} numbers, we extract, at \textbf{the} moment, 15, 000 million tons of \textbf{fossil} fuels every year. \textbf{The} International Energy Agency says that we will need about 43 million tons of critical materials in order to build out \textbf{the} clean energy economy.
So that ' s 300 times less stuff that is required.
It ' s worth saying that \textbf{the} impact of extracting 300 times less stuff upon nature is dramatically lower.
And I know that it ' s often argued that \textbf{the} \textbf{renewable} economy may be damaging for nature.
Again, this is a completely false statement.
It ' s in fact, our only chance of reducing \textbf{the} pressure upon our natural resources, because there ' s far less of it.
MZ : Of course, transitioning to using clean energy, upending industries, it has to happen across \textbf{the} world, not just here in \textbf{the} US.
And it can be hard to even imagine what life will be like.
But \textbf{carbon} scientist and futurist Julio Friedmann can help us.
Julio Friedmann : I think we need a compelling vision of \textbf{the} future that gets people excited about \textbf{the} energy transition and solving climate change.
Right now, it ' s like we ' re at a restaurant and there ' s two options : burnt toast, unflavored oatmeal.
And people are not excited about either, right? \textbf{The} burnt toast option for me is what I call \textbf{the} big switch.
Everything ' s \textbf{the} same, it ' s just clean, but it costs a lot of money and it ' s super hard to do. \textbf{The} other vision is kind of like,"Thou shalt not."That vision is a narrower sort of,"eat your peas,"hair shirt kind of vision of \textbf{the} future.
And a lot of people don ' t want that either.
And there ' s nothing in \textbf{the} physics, \textbf{chemistry} or \textbf{biology} that says we can ' t have a super interesting, exciting, vibrant world.
And that is compelling in a way that we have not heard yet.
MZ : Help me out here, what does that even look like?
JF : So, for lack of a better word, like, I like \textbf{the} Marvel cinematic \textbf{universe}.
MZ : OK, I ' ll go with you.
JF : So a place like Wakanda, right?
It ' s got flying cars, it ' s got maglev trains, and it ' s a totally fabulous place to live.
Like, they ' ve got lots of food.
That is built on abundant, sustainable, cheap energy.
MZ : OK, so like, we ' re not Wakanda, we ' re not in \textbf{the} Marvel \textbf{universe}.
This is real life, sort of.
How do we make this possible?
JF : So what can you actually do here on Earth, constrained by economics and engineering and all \textbf{the} rest of it, right?
And \textbf{the} answer is we can do an incredible amount.
Abundant, available where you want it, when you want it.
Everybody should have energy, including \textbf{the} three billion people who use less electricity than my refrigerator uses.
Well, \textbf{the} good news is, every day, \textbf{the} Earth receives 163, 000 terawatts of energy from \textbf{the} sun.
About half of that bounces back to space, but about 80, 000 terawatts arrive at \textbf{the} Earth in a form we can use.
For reference, today, \textbf{the} world uses about 26 terawatts of energy, and we got more than solar and wind.
We have geothermal, we have hydro, we have nuclear.
There ' s other kinds of clean energies.
And some of \textbf{the} best resources are, in fact, in \textbf{the} Global South.
These places are not simply future climate victims.
These places are latent energy superpowers.
My favorite example of this is Chile.
Almost eight years ago now their government said, wait a second, we can make \textbf{the} cheapest green electricity on Earth.
We have hydropower, we have solar power, we have wind power and \textbf{the} best resources anywhere in \textbf{the} world.
So let ' s start by making a lot of green electricity.
And then they said, because of that, we can also make green hydrogen cheaper than anywhere else.
And then Japan was like, hello, we want your green hydrogen.
Can you ship it to us?
MZ : Is that possible?
JF : Totally.
And Japan is like, we will also give you money to pay for these projects.
We ' ll loan you money to pay for these projects.
Could you please use Japanese technology?
So they ' re using Japanese \textbf{turbines}, electrolyzers, they ' re using Japanese ships to haul \textbf{the} ammonia around.
So Japan gets rich on this, Chile gets rich on this.
MZ : Win-win.
JF : Totally win-win.
MZ : Win-win-win, \textbf{planet}, too.
JF : And that model, other countries are looking at that, going, I want some of that.
And that ' s why I like this vision of abundance. 15 years ago, solar was \textbf{the} most expensive power.
Now it is \textbf{the} cheapest form of electricity, right?
Which is astonishing, it ' s totally amazing.
And a whole bunch of people are like, now we ' re done.
And I ' m like, no, we can do that with everything.
We can do that with wind.
We can do that with nuclear, we can do that with geothermal.
But these projects take time, money and people.
We need to develop \textbf{the} human capital as part of these investment projects for decades.
And innovation, we need much more energy in many places, much cheaper.
That ' s an innovation agenda.
For solar, that was \textbf{the} United States, Germany, China and others acting together.
Well, if that ' s \textbf{the} recipe, we can do that again.
We ' re already doing it with electric vehicles and clean hydrogen production.
We can certainly do it by turning electricity into fuels.
MZ : It sounds like \textbf{the} climate discussion needs to be kind of like improv, that it ' s a"yes and"sort of thing.
JF : I ' m so glad you said that.
So many people don ' t understand my love of improv comedy.
Yes, it is a"yes and"circumstance.
Because mostly it is about room for agreement.
Where do you find \textbf{the} point where everybody is like, I can agree on that?
One of \textbf{the} things I love about \textbf{the} infrastructure bill, one of \textbf{the} things I love about \textbf{the} Inflation Reduction Act, neither of them had \textbf{the} word climate in it.
They were both massive climate bills, right?
That was a way to get everybody on \textbf{the} same page.
Nationwide, they basically tripled \textbf{the} clean energy research budget.
That ' s a good start.
But that made it all of one percent of GDP.
And for comparison, we spend twice that on pharmaceuticals.
So I think it ' s fair to say we should get going.
We should spend a lot more money.
We ' ve done stuff like this before.
We ' ve done World War II, we ' ve done \textbf{the} Marshall Plan, like, we ' ve done \textbf{the} space shot, like we actually know how to move really quickly as a nation.
We did this for COVID.
Collective action, building together is what makes \textbf{the} difficult possible and nourishes \textbf{the} soul through mission and purpose.
We know what to do and we can act, not out of anger or fear, but out of generosity and common purpose, bringing aspiration and humility together.
We ' re going to build a thriving, vibrant, exciting world full of potential that ' s going to be built on \textbf{the} back of infrastructure, innovation and investment that will harness \textbf{the} abundant, sustainable, cheap energy that is our \textbf{planet} ' s endowment.
Thank you. [ Chapter 4 : A Just Transition ] FD : There is a high risk that when those exponential kick in, we see businesses running after those opportunities.
But then \textbf{the} distribution of \textbf{the} benefits being very unequal.
EZ : Many people like discussing \textbf{the} importance of \textbf{the} just energy transition.
And a lot of people don ' t understand what justice really means.
If we are to reach a revolution, a \textbf{renewable} energy future has to be inclusive.
AC : We have this push for what we ' re calling a just and equitable energy transition.
We need to make sure that we are not exacerbating inequality, but rather we are ensuring that that gap is closed and that \textbf{the} groups that have been historically marginalized are not once again being marginalized in a new way.
NT : Most of \textbf{emissions} come from wealthy, energy-guzzling citizens, but a lot of \textbf{the} economic benefit and human development benefit will come not from decarbonizing high-carbon lives and lifestyles, but from providing clean, distributed energy and tools and solutions to those who don ' t have access to any of them.
LN : When we speak of countries in need of economic development, we mostly speak of countries that are hit by \textbf{the} climate crisis \textbf{the} hardest right now.
It should be \textbf{the} biggest wish and job and task for all those countries who have burned so much \textbf{fossil} fuel already to support any other country in \textbf{the} so-called Global South, to get an economy running without \textbf{fossil} fuels, to not repeat all \textbf{the} mistakes that we have made.
HD : To ensure that \textbf{the} solar revolution that has been seen around \textbf{the} world comes to places like Africa, we need to be really conscious that \textbf{the} history and \textbf{the} lived experience when it comes to energy is very different for poor countries.
Rebekah Shirley : We love to remind Global South communities, especially Africa, about \textbf{the} vast \textbf{renewable} energy potential that they have, which is true.
But \textbf{the} finance flows to deliver on that potential remain very scarce.
Africa, despite being 90 percent of those across \textbf{the} entire globe that still do not have access to electricity and being 17 percent of \textbf{the} global population, we ' re actually two percent of international finance.
Money doesn ' t actually flow to these spaces naturally, and we need to do a lot of work to get global finance to these spaces.
MZ : A world of abundant clean energy for everyone.
That is \textbf{the} goal.
But what about in \textbf{the} short term?
Take Detroit.
This once-wealthy city has been through some seriously hard times.
How is it rebuilding itself so that everyone here can thrive while dealing with extreme heat, flooding, and other climate problems?
Anika Goss is looking to Detroit ' s past to plan its future.
AG : I am a third-generation Detroiter.
My grandmother moved to Detroit in 1936 during \textbf{the} Great Migration and brought all of her southern ways with her.
She owned a home and knew that home ownership would create wealth and opportunity for her growing family.
Up until \textbf{the} late 1950s, Detroit was a haven for middle-class families living in neighborhoods where there was green space and community connectivity and opportunity.
My grandmother, she had this amazing garden, and it was just abundant.
There were all kinds of flowers and food.
You know, she grew vegetables back there, and it was just a magical place for a little girl.
MZ : I mean, what you ' re describing is a Detroit that, yes, has huge problems, segregation.
But at \textbf{the} same time, a woman, she can have a job, she can own a house, she can raise her family, she can have fresh fruits and vegetables.
AG : It was certainly a place where you could gain economic opportunity faster than you can now.
My grandmother ' s Detroit is not \textbf{the} Detroit that I live in today.
All of those sites where there was manufacturing and industrial sites that led to our economic boom, many stand vacant and abandoned.
These industrial sites have led to dangerous contamination to our land and our water and our air.
MZ : So, just lay it out for us.
For people who maybe don ' t see, like, how is there a link between economic inequality and climate change?
AG : You know, I ' m so surprised that it feels like separate problems.
But I ' m also really glad that people are asking.
Because if you ' re living around and working in a community that is low-income, where there are fewer jobs, \textbf{the} housing is deteriorated, this is a neighborhood that ' s also more likely to be closer to a factory that ' s emitting pollutants, likely to have water that ' s contaminated, probably has fewer parks and fewer intentional green spaces. \textbf{The} people who are living in those communities are at risk for climate impacts first.
So we have to do something in those neighborhoods right away.
MZ : So what can this healthier future look like?
Anika showed me some of \textbf{the} work she ' s done in partnership with different neighborhoods in Detroit, to make them both more enjoyable and more resilient.
First, she took me to \textbf{the} Circle Forest.
It ' s an oasis right in \textbf{the} middle of \textbf{the} city.
AG : This is six lots that were vacant, and we all worked together with this community to build a forest.
MZ : Beautiful and a great example of how nonprofits can work with communities to transform them into sustainable, welcoming places.
AG : We might have a good idea about gardens and tree planting, but this is where people live.
And so you ask first.
This is a new effort, of this idea of a resilient neighborhood or calling it that.
So these are large tracts of vacancy that can be used to provide environmental \textbf{sustainability}, community resilience and beauty.
Watch for \textbf{the} poison ivy.
It ' s not pretend.
MZ : Next we met up with Katrina Watkins, \textbf{the} founder and CEO of \textbf{the} Bailey Park Neighborhood Development Corporation, a group revitalizing this historic area by turning vacant lots into parks that can stand up to extreme weather.
Katrina Watkins : My family has been in this neighborhood since 1946.
I grew up hearing \textbf{the} stories from my dad, him telling me about how people thought it was a ghetto, but it wasn ' t.
Those were families, and people worked and had businesses, and people looked out for each other.
And it was those type of stories that made me say, wow, we just really need to bring that back.
MZ : So then briefly fill us in about how things changed.
KW : Things deteriorated in this neighborhood. \textbf{The} homes got torn down, left just a lot of vacant lots, a lot of overgrowth.
It wasn ' t safe.
I would see my little nieces and nephews.
They would just be on their little skateboards, you know, not paying attention, \textbf{the} stop sign is covered.
And me and my dad was like, we just have to do something.
So we started out as Bailey Park Project, focusing on \textbf{the} park and green infrastructure and trees and flowers.
And then over \textbf{the} course of \textbf{the} years, we started doing more community development.
MZ : I wonder, like, how do you describe it to people when you explain what ' s happening, what you ' re trying to do for your neighborhood, \textbf{the} changes that you ' re making and how it relates to climate change?
KW : One really great example, when people come to \textbf{the} park, like right now, this is like, literally a heat \textbf{island} because our trees haven ' t grown.
So it is very hot.
So sometimes it ' s so hot that \textbf{the} kids don ' t come out and play until \textbf{the} evening.
That education of knowing \textbf{the} benefit of trees and cooling down an environment, how it makes \textbf{the} air cleaner.
There ' s a lot of people that suffer from asthma.
So trees can make a difference.
And if you go to places like Grosse Pointe, you will see this beautiful tree canopy down.
Like their homes are probably 20 degrees cooler you know, than our homes, AG : Having a lot of green space actually contributes to \textbf{the} climate resiliency for this neighborhood.
So all of this grass actually can manage stormwater better. \textbf{The} idea of leading new development in this neighborhood with green, what it says to me is that \textbf{the} residents are really envisioning \textbf{the} neighborhood that they want to live in.
And that I just -- I ' m really excited about that.
MZ : Turning great loss and upheaval into even greater opportunities.
These Detroiters are creating a blueprint for an urban future that every city can learn from. [ Chapter 5 : Agriculture ] TED-Ed : About 10, 000 years ago, humans began to farm.
This agricultural revolution was a turning point in our history that enabled people to settle, build and create.
In short, \textbf{agriculture} enabled \textbf{the} existence of civilization.
Today, approximately 40 percent of our \textbf{planet} is farmland.
Spread all over \textbf{the} world, these agricultural lands are \textbf{the} pieces to a global puzzle we are all facing.
In \textbf{the} future, how can we feed every member of a growing population a healthy diet?
Meeting this goal will require nothing short of a second agricultural revolution. \textbf{The} first agricultural revolution was characterized by expansion and exploitation, feeding people at \textbf{the} expense of forests, wildlife and water, and destabilizing \textbf{the} climate in \textbf{the} process.
That ' s not an option \textbf{the} next time around.
Agriculture depends on a stable climate with predictable seasons and weather patterns.
This means we can ' t keep expanding our agricultural lands, because doing so will undermine \textbf{the} environmental conditions that make \textbf{agriculture} possible in \textbf{the} first place.
Instead, \textbf{the} next agricultural revolution will have to increase \textbf{the} output of our existing farmland for \textbf{the} long term while protecting biodiversity, conserving water, and reducing pollution and \textbf{greenhouse} gas \textbf{emissions}.
So what will \textbf{the} future farms look like?
In Bangladesh, Cambodia and Nepal, new approaches to rice production may dramatically decrease \textbf{greenhouse} gas \textbf{emissions} in \textbf{the} future.
Rice is a staple food for three billion people and \textbf{the} main source of livelihood for millions of households.
More than 90 percent of rice is grown in flooded paddies, which use a lot of water and release 11 percent of annual methane \textbf{emissions}, which accounts for one to two percent of total annual \textbf{greenhouse} gas \textbf{emissions} globally.
By experimenting with new strains of rice, irrigating less and adopting less labor-intensive ways of planting seeds, farmers in these countries have already increased their incomes and crop yields while cutting down on \textbf{greenhouse} gas \textbf{emissions}.
MZ : Related technologies are also taking root in \textbf{the} US.
Father and daughter farmers, Jim and Jessica Whitaker, have come up with their own strategies to reduce methane that are working on their farm in Arkansas.
But convincing everyone else to change is their next big challenge.
So let ' s go back to how you got to be such a big supplier of rice in this country.
You tell a story about how you were 22 and you were like, I ' m going to start my own farm, but -- not that easy.
Jim Whitaker : No, it ' s not.
Let me tell you something about renting a farm when you ' re 22 years old.
No one rents you a farm unless no one else wants to rent it.
It was a big piece of land, it was cash rent, I mean, \textbf{the} stakes were set.
We were doomed to fail.
And we weren ' t focused on environmental \textbf{sustainability} back then.
We were focused on economic \textbf{sustainability}.
How do we make higher yield?
How do we use less fertilizer?
How do we get to \textbf{the} next year and feed our family?
It happens that economic and environmental \textbf{sustainability} go hand in hand along with social \textbf{sustainability}.
As we used less fertilizer, used less water, our yield started going up.
MZ : Can you spell out for me what \textbf{the} problem is with rice when it comes to \textbf{the} environment?
JW : Soil is full of microbes.
And what they ' re doing in \textbf{the} soil on a basic level is they ' re down there chewing away, eating up \textbf{the} biomass that ' s left over from \textbf{the} crop before.
And out of that comes methane, \textbf{the} same way it happens in cattle operations or whatever.
MZ : But you decided to do it a different way.
JW : We do it a little bit differently.
Arkansas is a very water-rich state, but we have a dry season, right in \textbf{the} middle of summer is a dry season.
So when we dry \textbf{the} soil, that anaerobic microbial that ' s living in that \textbf{soup}, that mushy soil, either dies or goes dormant.
So we break his life cycle.
MZ : Oh!
JW : And then it stops emitting methane.
But it doesn ' t hurt \textbf{the} rice.
That actually helps rice, when you use less water.
We have less runoff, less erosion, less nutrients leaving our field.
Nothing leaves our field unless we want it to.
Jessica Whitaker Allen : We started tracking this.
We saw that our water use was down and we thought, OK, well, we need to add to that, let ' s add more data points to that.
Let ' s add our fertilizer, let ' s add our fuel usage.
And so we really started recording all this stuff by hand in just an Excel spreadsheet.
We didn ' t know what we were doing.
MZ : But you knew you were on to something.
JWA : Yes, so we started recording everything.
MZ : Do you remember if there was a moment when \textbf{the} two of you thought, like, what we ' re doing, it shouldn ' t be special to us.
Why aren ' t more people -- JW : To collect \textbf{the} data, have \textbf{the} equipment and \textbf{the} cloud-based technology and \textbf{the} stuff in \textbf{the} field to record it, it ' s just massively time-consuming, and consumers don ' t pay for it.
JWA : Well and people like us, why would you want to do it?
I mean, take a walk through your \textbf{supermarket}.
Everything ' s sustainable, everything ' s greenwashed.
You can put whatever you want to on a package, there ' s no rules, there ' s no standards.
So we ' re trying to make that standard.
We live and work in southeast Arkansas.
Our town has 4, 000 residents and one stoplight. \textbf{The} nearest airport, Starbucks shopping mall, Whole Foods, is two hours away in any direction.
Without places like this and farmers like us, you ' d be hungry, naked and sober.
I ' m a farmer, but not in \textbf{the} way you might think.
I don ' t drive \textbf{the} tractors, and I don ' t plant \textbf{the} rice.
But I know how to take what my dad has learned and what he is doing and share that with others.
We are going to work with farmers in southeast Arkansas to educate them about \textbf{the} benefits of growing sustainable rice.
JW : Let me tell you why that ' s important.
There are 400 million acres of rice grown globally.
And our methods, if used, can reduce \textbf{greenhouse} gas by 50 percent, reduce water use by 50 percent, increase yields to feed a hungry world.
MZ : We hear a lot about how weather has gotten more extreme, flooding, heat.
What ' s it been like in Arkansas?
JW : So in 2016, we had a 500-year flood.
We were...
We get a picture of our place flooded. 14 \textbf{inches} of rain.
That was a 500-year flood.
Hurt a lot of people.
And then we had a 1, 000-year flood.
And we lost so many acres.
And you know, when houses get flooded in rural communities, they don ' t come back.
I mean, when you have whole neighborhoods get flooded, they don ' t rebuild them, they don ' t come back.
JWA : I feel like almost every day I ' m getting a notification from \textbf{the} weather Channel of an excessive heat warning or thunderstorm warning or hail threat.
You know, you don ' t really know what you ' re going to expect.
MZ : When you think about what it will be like, I don ' t know, five years, 10 years, for someone to go into \textbf{the} \textbf{supermarket} and think, oh, right, I need to get rice off my grocery list.
What ideally would you like to see on \textbf{the} shelf there?
JW : So I think it ' s going to change, because I like to study what companies want.
I mean, if you figure out what they want, then you know what to give them.
And when you read their \textbf{sustainability} goals, they ' re all saying they ' re going to be net zero. \textbf{The} only way they can be net zero is they bring \textbf{the} farmer into \textbf{the} loop.
So if we ' re able to do this, in five years, everybody ' s doing these practices, \textbf{the} US rice industry will be \textbf{the} most sustainable industry in \textbf{the} world.
MZ : Can you do it?
JW : Oh, yeah.
It ' s about to happen.
MZ : You ' ve just heard a lot of ideas, a lot of stories and a lot of big plans.
So we want to wrap things up by asking our experts : What if we actually achieved our climate goals?
What would that future even look like?
They, of course, had a lot of thoughts.
We ' ll let you pick your favorites.
Thanks for watching.
NT : If we are successful in addressing \textbf{the} climate crisis, we massively reduce \textbf{the} vulnerability and improve \textbf{the} economic and \textbf{the} creative prospects of billions of people around \textbf{the} world who currently are basically excluded from all of \textbf{the} opportunities which we have, which come from having abundant material and energy all around us.
That ' s \textbf{the} real prize.
FD : A world where we have had that exponential would be a world in which we have clean, cheap energy for all.
KB : At \textbf{the} end of this decade, \textbf{the} International Energy Agency believes that we will have enough capacity to be producing 10 billion, 10 billion batteries a year.
It ' s enough for every person on \textbf{the} \textbf{planet} to have their own battery.
AC : I ' m constantly surrounded by young people just like me who are really, really, really passionate about \textbf{the} world that we live in.
For obvious reasons.
It ' s \textbf{the} world that we will inherit.
EZ : Women are very important for \textbf{the} energy transition.
So if we ' re talking about a world where we have, you know, where we tackle climate change, all I want to see is, you know, women in power.
I want to see more solutions being driven by young people.
And I want our politicians to be truly and meaningfully engaged in making change.
NG : We always tend to underestimate \textbf{the} pace of ultimate adoption and change.
Today, there are more than 150 pathways, scientific pathways that we could take to limit global warming.
Al Gore : If we get to true net zero, astonishingly, global temperatures will stop going up with a lag time of as little as three to five years.
Colin Averill : We ' ve been able to accelerate tree growth and \textbf{carbon} capture and tree stems by 30 to 70 percent.
Dan Jørgensen : Offshore wind has \textbf{the} potential to cover \textbf{the} current electricity demand of \textbf{the} entire world not once, not twice, 18 times.
NG : We need to decarbonize \textbf{the} old economy that we have and invest to create \textbf{the} new economy that we need Paul Hawken : Regenerative farming is fairly simple in concept.
It means creating \textbf{the} conditions for more life.
Josephine Phillips : We need to buy less stuff, and we need to look after what we buy.
Valuing \textbf{the} things that we own is a climate solution.
Vishaan Chakrabarti : We can all live in this transit-rich, carbon-negative, affordable way and leave \textbf{the} vast majority of \textbf{the} \textbf{planet} for nature, for \textbf{agriculture}, for clean oceans.
We can do this.
Susan Ruffo : \textbf{The} ocean is a powerful source of solutions that we ' ve overlooked for far too long.
Al Roker : You have to be engaged.
Let your elected officials know that this is important to you.
You have to vote, you have to go out there and support politicians who are going to support our \textbf{planet}.
Peggy Shepard : We can create a legacy of environmental quality and climate resilience for all.
Fehinti Balogun : You are not powerless that \textbf{the}"oh, there ' s nothing we can do, just keep going"is symptom, not cure.
But you are needed.
Avinash Persaud : \textbf{The} power of what we can do when it matters to us is unlimited.
CF : It ' s not that we ' re simply going to avoid \textbf{the} worst impacts of climate change.
It ' s that we ' re going to improve \textbf{the} quality of life for humans, both in \textbf{the} urban and in \textbf{the} rural sector, and for non-humans.

% matched lemmas: agriculture, backwards, biology, carbon, chemistry, coal, electrical, emission, fossil, greenhouse, inch, island, occur, oil, planet, renewable, shipping, smartphone, soup, spring, supermarket, sustainability, the, transportation, turbine, universe
\end{textsample}
